
\section{死锁概念}
\concept{死锁定义}
\begin{points}
  \item 死锁是多个进程因竞争系统资源而造成的一种僵局，若无外力作用，这些进程都将永远不能向前推进。
  \item 若进程集合中的每个进程都在等待只能由该集合中其他进程引发的事件，则该集合发生死锁。
  \item 典型例子：P1 占有打印机等待读卡机，P2 占有读卡机等待打印机，双方都无法继续。
\end{points}
\plain{死锁就是大家都拿着别人需要的资源不放，结果谁都走不下去。}
\exam{常考四个必要条件、预防/避免/检测解除区别、银行家算法安全序列。}


\concept{死锁与饥饿、死循环}
\begin{points}
  \item 死锁：一组进程相互等待，若无外力永远不能推进。
  \item 饥饿：某进程长期得不到所需资源，但系统整体仍在推进。
  \item 死循环：一个进程自身逻辑错误导致反复执行，不一定涉及资源等待。
  \item 活锁：进程没有阻塞，但不断改变状态、相互让步，导致没有实际进展。
\end{points}
\plain{死锁就是大家都拿着别人需要的资源不放，结果谁都走不下去。}
\exam{常考四个必要条件、预防/避免/检测解除区别、银行家算法安全序列。}


\section{资源分类与死锁必要条件}
\concept{资源分类}
\begin{points}
  \item 可剥夺资源：可从进程手中强行取回且不破坏计算，例如 CPU、内存页框。
  \item 不可剥夺资源：一旦分配，不能强行剥夺，否则会造成错误，例如打印机、部分外设。
  \item 可重用资源：使用完可再次分配，如设备、内存。
  \item 消耗性资源：使用后消失，如消息、信号。
  \item 死锁主要与不可剥夺的可重用资源有关。
\end{points}
\plain{把这个知识点放回本章主线理解：它回答的是 OS 为了管理资源、隐藏硬件、保证并发正确性而采用的某个对象、规则或步骤。}
\exam{常考选择/填空/简答，让你判断概念归属、比较相近概念，或按“定义-作用-机制-易错点”展开。}


\concept{死锁产生的四个必要条件}
\begin{points}
  \item \key{互斥条件}：资源一次只能被一个进程使用。
  \item \key{请求和保持条件}：进程已保持至少一个资源，同时又请求新的资源。
  \item \key{不可剥夺条件}：资源不能被强行夺走，只能由占有进程主动释放。
  \item \key{循环等待条件}：存在进程资源等待环，$P_1$ 等 $P_2$，$P_2$ 等 $P_3$，$\ldots$，$P_n$ 等 $P_1$。
\end{points}
\plain{死锁就是大家都拿着别人需要的资源不放，结果谁都走不下去。}
\exam{常考四个必要条件、预防/避免/检测解除区别、银行家算法安全序列。}

\section{处理死锁的基本方法}
\concept{四类方法总览}
\begin{points}
  \item 死锁预防：通过制度性限制破坏死锁必要条件之一。
  \item 死锁避免：运行时根据资源需求判断分配是否安全，典型算法是银行家算法。
  \item 死锁检测与解除：允许死锁发生，定期检测，发生后恢复。
  \item 鸵鸟策略：忽略死锁问题，适用于死锁概率很低且处理成本过高的系统。
\end{points}
\plain{把这个知识点放回本章主线理解：它回答的是 OS 为了管理资源、隐藏硬件、保证并发正确性而采用的某个对象、规则或步骤。}
\exam{常考选择/填空/简答，让你判断概念归属、比较相近概念，或按“定义-作用-机制-易错点”展开。}


\concept{死锁预防}
\begin{points}
  \item 破坏互斥条件：把独占设备改造成共享设备，例如 SPOOLing 把打印机虚拟成共享设备。但并非所有资源都能共享。
  \item 破坏请求和保持条件：进程运行前一次性申请所有资源，或申请新资源前释放已占资源。缺点是资源利用率低、可能饥饿。
  \item 破坏不可剥夺条件：资源申请失败时剥夺其已占资源。适用于 CPU、内存等可剥夺资源，不适合打印机等。
  \item 破坏循环等待条件：对资源统一编号，进程按编号递增顺序申请资源。
\end{points}
\plain{死锁就是大家都拿着别人需要的资源不放，结果谁都走不下去。}
\exam{常考四个必要条件、预防/避免/检测解除区别、银行家算法安全序列。}


\concept{安全状态与不安全状态}
\begin{points}
  \item 安全状态指系统能找到一个安全序列，使所有进程都能按某种顺序获得所需资源并完成。
  \item 不安全状态不一定已经死锁，但可能发展为死锁。
  \item 死锁状态一定是不安全状态；不安全状态不一定是死锁状态。
\end{points}
\plain{进程不是一直运行，它会在就绪、运行、阻塞等状态之间因调度或等待事件而移动。}
\exam{常考三/五/七状态转换图，尤其运行到阻塞、阻塞到就绪、就绪到运行的触发条件。}


\concept{银行家算法}
\begin{points}
  \item 银行家算法用于死锁避免，假设系统知道每个进程的最大资源需求。
  \item 关键数据结构：Available 可用资源，Max 最大需求，Allocation 已分配，Need 尚需，且 $Need=Max-Allocation$。
  \item 资源请求满足 $Request_i\le Need_i$ 且 $Request_i\le Available$ 后，先试探性分配。
  \item 再执行安全性检查：寻找一个进程满足 $Need_i\le Work$，假设其完成并释放资源，更新 $Work=Work+Allocation_i$，直到所有进程完成或找不到下一个。
  \item 若试分配后仍安全，则真正分配；否则让进程等待。
\end{points}
\plain{死锁就是大家都拿着别人需要的资源不放，结果谁都走不下去。}
\exam{常考四个必要条件、预防/避免/检测解除区别、银行家算法安全序列。}


\concept{资源分配图}
\begin{points}
  \item 资源分配图用圆表示进程，用方框表示资源类型，方框内点表示资源实例。
  \item 请求边：进程 $\rightarrow$ 资源；分配边：资源实例 $\rightarrow$ 进程。
  \item 若每类资源只有一个实例，图中有环等价于发生死锁。
  \item 若资源类型有多个实例，有环只是死锁的必要条件，不一定死锁。
\end{points}
\plain{死锁就是大家都拿着别人需要的资源不放，结果谁都走不下去。}
\exam{常考四个必要条件、预防/避免/检测解除区别、银行家算法安全序列。}


\concept{死锁检测与解除}
\begin{points}
  \item 死锁检测允许资源正常分配，周期性检查是否存在死锁。
  \item 检测后可通过撤销进程、回滚进程、剥夺资源等方式解除死锁。
  \item 选择牺牲进程时可考虑优先级、已运行时间、已占资源、还需资源、回滚代价等。
  \item 缺点是检测和恢复都有代价，且可能造成进程饥饿。
\end{points}
\plain{死锁就是大家都拿着别人需要的资源不放，结果谁都走不下去。}
\exam{常考四个必要条件、预防/避免/检测解除区别、银行家算法安全序列。}


